\section{Problem Formulation}
An observation is an invoice-level set of distinct product codes. Quantity is used to identify positive sales records, but repeated units of a product do not increase its basket membership. Let $D_t$ be the set of eligible orders in month $t$, with $N_t=|D_t|$. A rule $A\rightarrow B$ has a nonempty antecedent $A$ and a single-product consequent $B$, with $A\cap B=\varnothing$. The combined itemset contains at most three products. The arrow denotes conditional co-occurrence within an order, not purchase sequence or causation.

Let $n_t(X)$ count orders containing every product in itemset $X$. We calculate
\begin{align}
\operatorname{support}_t(A\rightarrow B) &= \frac{n_t(A\cup B)}{N_t},\\
\operatorname{confidence}_t(A\rightarrow B) &= \frac{n_t(A\cup B)}{n_t(A)},\\
\operatorname{lift}_t(A\rightarrow B) &= \frac{n_t(A\cup B)N_t}{n_t(A)n_t(B)}.
\end{align}
Confidence describes the fraction of antecedent-containing orders also containing the consequent. Lift adjusts that fraction for the consequent's prevalence; a lift above one indicates positive co-occurrence relative to independence. Neither metric is a causal estimate.

\subsection{A fixed definition of retention}
All discovery configurations use the same monthly evaluation criteria:
\begin{center}
\begin{tabular}{lr}
\toprule
Requirement & Monthly threshold\\
\midrule
Support & $\geq 0.005$\\
Confidence & $\geq 0.50$\\
Lift & $\geq 1.20$\\
Antecedent count & $\geq 20$ orders\\
\bottomrule
\end{tabular}
\end{center}
The thresholds are operational choices for this project, not universal statistical standards. The antecedent-count floor limits conclusions based on very few relevant orders. A rule is marked evaluable when its antecedent occurs at least 20 times and its consequent appears in the month. If the antecedent is absent, confidence is undefined; if either marginal is zero, lift is undefined. Such cases remain in the frozen cohort and are recorded separately rather than silently removed or converted to numerical zero.

For frozen discovery cohort $R$, let $I_t(r)$ equal one when all monthly requirements hold. Define
\begin{equation}
\operatorname{Retention}_t=\frac{\sum_{r\in R}I_t(r)}{|R|},\qquad
\operatorname{Coverage}_t=\frac{\#\{r\in R:r\text{ is evaluable in }t\}}{|R|}.
\end{equation}
Conditional retention uses the evaluable subset as its denominator and is reported separately. The primary retention measure therefore captures both adequate observation and criterion satisfaction; it is not a pure estimate of behavioral change. A rule's persistence score is the number of held-out months in which it is retained, from zero to four. The mean monthly rate weights the four months equally; here it also equals the retained fraction of all rule-months because the cohort size is constant.
For nonzero coverage, these definitions give the exact decomposition
\begin{equation}
\operatorname{Retention}_t
=\operatorname{Coverage}_t\,
\operatorname{ConditionalRetention}_t.
\label{eq:decomposition}
\end{equation}
This identity also holds for pooled counts, but not generally for the product of unweighted monthly means. It separates an evidence-availability problem from failure of the observed association to meet the declared requirements.

\subsection{Consequences of finite monthly counts}
The support and confidence gates imply the integer-count requirement
\begin{equation}
n_t(A\cup B)\geq
\max\!\left\{\left\lceil0.005N_t\right\rceil,\,
\left\lceil0.50n_t(A)\right\rceil\right\}.
\end{equation}
The lift and evaluability gates must also hold. In August, support alone requires at least six joint orders; in November it requires 13. At the minimum antecedent count of 20, confidence requires at least ten joint orders. The antecedent floor itself corresponds to 1.67\% of August orders but 0.81\% of November orders. These are consequences of the same policy applied at different monthly volumes, not month-specific tuning. A rule frequent over six pooled months need not supply adequate evidence in each later month, so low coverage alone does not establish behavioral drift.

